\begin{comment}
=== AI-INSTRUCTIONS (NEVER DELETE OR ALTER THIS BLOCK) ===
1. LHS STYLE: Explanations on LHS use \footnotesize.
2. SPATIAL LIMITS: Keep the minted code to ~40 lines maximum.
3. REVIEW NOTES:
   - [Tim Note]: Pending detailed review.
==========================================================
\end{comment}

\begin{frame}[fragile]{Base Primitives: Columns \& Summaries}
\begin{columns}[T]
  \begin{column}{0.4\textwidth}
    \footnotesize 
    \textbf{Lines 86-110:}\\
    This chunk handles base data compression. 

    \vspace{0.2cm}
    \rc{1} \textbf{Sym:} Tracks counts of symbols.
    
    \vspace{0.15cm}
    \rc{2} \textbf{Num:} Uses Welford's algorithm to track mean and variance on the fly.
  \end{column}
  
  \begin{column}{0.6\textwidth}
\begin{minted}{python}
# ---- 1. Columns ----
class Sym: |\rc{1}|
  def __init__(i, txt="", at=0):
    i.txt, i.at, i.n, i.has = txt, at, 0, {}
  def add(i, x):
    if x != "?":
      i.n += 1
      i.has[x] = 1 + i.has.get(x, 0)
  def mid(i): 
    return max(i.has, key=i.has.get) if i.has else "?"

class Num: |\rc{2}|
  def __init__(i, txt="", at=0):
    i.txt, i.at, i.n = txt, at, 0
    i.mu, i.m2 = 0, 0
    i.lo, i.hi = 1E32, -1E32
  def add(i, x):
    if x != "?":
      i.n += 1
      d = x - i.mu
      i.mu += d / i.n
      i.m2 += d * (x - i.mu)
      i.lo = min(x, i.lo)
      i.hi = max(x, i.hi)
  def div(i): 
    return (i.m2 / (i.n - 1))**.5 if i.n > 1 else 0
\end{minted}
  \end{column}
\end{columns}
\end{frame}
